\section{Related Work}
\label{sec:related}

\subsection{CCNet and Perplexity-Based Quality Filtering}

\citet{Wenzek2019CCNet} introduced CCNet, a pipeline for extracting high-quality monolingual text from CommonCrawl by training language-specific KenLM $n$-gram models on Wikipedia and scoring web documents by perplexity. The original design is explicitly per-language: each language receives its own KenLM model trained on its Wikipedia, and cutoffs are applied from language-specific tercile boundaries stored in \texttt{cutoff.csv}. This per-language calibration was not an afterthought --- it acknowledges that perplexity scales depend on the reference corpus and should not be compared across languages.

However, practitioners adapting CCNet for large-scale multilingual pipelines have often simplified this design by applying a single global percentile threshold to the combined multi-language document pool. Dolma \cite{Soldaini2024Dolma} applies Pythia-based perplexity scores; RedPajama-V2 \cite{Weber2024RedPajama} provides pre-computed \texttt{ccnet\_perplexity} as a quality signal metadata field intended for downstream filtering. In each case, the path from pre-computed perplexity to actual data selection --- including whether filtering is applied globally or per-language --- is left to the practitioner. Our work shows that the most natural choice (global $k$-th percentile) has consequences that are not obvious from the CCNet paper's description.

\subsection{Multilingual Dataset Quality Audits}

\citet{Caswell2021Quality} provide the most systematic audit of per-language quality differences in multilingual web corpora, evaluating over 100 languages in mC4. They document severe quality disparities in low-resource languages and argue that per-language quality evaluation is necessary --- but they do not measure the associativity between language membership and retention rate under a global threshold. Their findings are qualitative and concern different filtering signals (not \texttt{ccnet\_perplexity} specifically).

\citet{Adelani2023Better} audit quality filtering for African languages and find that global filters consistently fail for low-resource languages. \citet{Niyomugabo2025BhashaKritika} (BhashaKritika) identify that per-language KenLM scoring is necessary even as a prerequisite quality evaluation step. These works motivate per-language approaches but do not provide quantitative $V$ measurements for a specific dataset-signal-threshold combination that practitioners can calibrate correction strategies against.

\subsection{Retention Rate Tuning for Multilingual Filtering}

The work closest to ours is \citet{Turki2026CrossLingual}, who study cross-lingual quality classifiers and find that retention rate tuning is necessary when applying multilingual quality filters --- globally calibrated classifiers systematically under-retain some language groups. They adopt percentile-based (relative) thresholds computed per regression head as a best practice. Importantly, they endorse per-language percentile calibration as superior to absolute threshold calibration, providing indirect prior support for the correction hypothesis we formulate. However, their setting differs: they work with neural classifiers trained on the mC4 corpus, not with pre-computed KenLM perplexity signals from RedPajama-V2.

\citet{Jansen2022Perplexed} find that standard perplexity filtering breaks down on multilingual heterogeneous data. \citet{Singh2024Repetition} find that per-language filtering strategies produce qualitatively different diversity-quality tradeoffs than English-centric approaches.

\subsection{RedPajama-V2 and Quality Signals}

\citet{Weber2024RedPajama} describe the RedPajama-V2 dataset and its quality signal pipeline. They acknowledge that ``ML-based quality signals have been reported to lead to biases or underrepresent minorities'' and include \texttt{ccnet\_perplexity} as one of 40+ quality signals computed for 5 languages (en, de, fr, es, it). They do not quantify the language-group retention disparity produced by global thresholding on any individual quality signal. Our work fills this gap for \texttt{ccnet\_perplexity}, providing a calibrated Cram{\'e}r's $V$ measurement that the RedPajama-V2 paper itself lacks.

\subsection{Our Position}

Our work is the first to quantify the language-group retention disparity (Cram{\'e}r's $V$) from global \texttt{ccnet\_perplexity} thresholding on RedPajama-V2, across multiple threshold levels, with Holm-corrected significance. This fills the gap between CCNet's per-language design intent \cite{Wenzek2019CCNet} and the practitioner-observed disparity \cite{Caswell2021Quality} by providing a precise, reproducible measurement. The $V = 0.40$--$0.57$ baseline we establish enables future correction studies to be properly powered and evaluated.
